\chapter{Planning}\label{chapter:planning}

% Description of what the plan is for this game engine.


My goal for this project is to start creating a game engine that I can work on in the future, and hopefully gather many of my friends to help me as well. 
Given the complexity and scope of this endeavour, it's crucial to start with a solid foundation.
For this regard, I have looked into many of the popular game engines as outlined in \autoref{chapter:research} and tried to pick the best ideas from each. 
For a successful open source project, it is important to make contributing as easy as possible, while also planning for growth of the project. Keeping this in mind was really important for me while choosing technologies to use, and will be detailing more in \autoref{}

\section{Modular Design}

Both Unity and Unreal engines have a modular design to ensure flexibility, extensibility, and manageability. Modules allow for separation of concerns, making it easier for contributors to understand, modify, and test individual components of the engine without interfering with others. 
A well-structured modular design allows for parallel development of different aspects of the game. This is really important for me as this large project is a collaborative effort from the start, with my friend making the renderer.
On to of this, modular design also enables the engine to adapt to a wide range of game genres and requirements, while not bloating the core of the engine with code that is not useful for the given game. For example, a VR module can only be enabled and loaded when it is actually necessary. Modularity is therefore a critical feature of the engine I want to create.

\section{Unified Editor}

The other important take away from examining existing game engines was that while many allowed for, most of them didn't need to compile the engine and editor for the given game. On top of this, another common aspect was that modern engines had all the editor tools integrated into a single application. This made development much easier and faster to learn. These mean a lot of time can be saved while developing the actual games, but they definitely add to the complexity of the project.

To give the engine the ability to load in the game code while running and not need to compile it together with the engine, the game code needs to be a separate compilation step. The engine needs to be able to load in this separate code and also reload it when it gets recompiled. This is a more tricky problem than it would first seam, as it needs to correctly link the external code.

To create a game engine where all necessary authoring tools are built into the editor is certainly possible, but would take exponentially longer than this thesis project allows. Because of this I will be not focusing on this aspect but leaving it for future development. In the meantime, I will be using external industry standard editors to create assets for testing.

\section{Entity system}

Every game engine needs a way to represent the game world in code in some way. In Unity there are the \verb|GameObject|s, in UE they are \verb|Actor|s, while valve's Source engine calls them Entities, but they fulfil a functionally similar role in making a game. As we can see from just these examples, there isn't an industry standard way of naming these building blocks, and their implementation also differs between engines greatly.

While in Unity and Unreal the user adds Components to create functionality, in Source all entities are pre-defined, and new ones can only be added using \Cpp \footnote{There are some Source games that support easier scripting with VScript \cite{valve-vscript}}. But Unity and UE also differ in their approach where Unity only has a flat list of components, UE allows for a whole internal hierarchy of components.